\uuid{vWk5}
\exo7id{601}
\titre{exo7 601}
\auteur{gourio}
\organisation{exo7}
\datecreate{2001-09-01}
\isIndication{false}
\isCorrection{false}
\chapitre{Suite}
\sousChapitre{Suite récurrente linéaire}
\cpge{mpsi-mp2i.nombres-reels-suites.c, pcsi.nombres-reels-suites.b, ptsi.nombres-reels-suites.b, mpsi-mp2i.nombres-reels-suites.d, pcsi.nombres-reels-suites.c, ptsi.nombres-reels-suites.c, mpsi-mp2i.nombres-reels-suites.i, pcsi.nombres-reels-suites.g, ptsi.nombres-reels-suites.g, mpsi-mp2i.analyse-asymptotique.c, pcsi.analyse-asymptotique.c, ptsi.analyse-asymptotique.c, bcpst1.suites-reelles, bcpst1.limites-continuite.b}
\module{Analyse}
\niveau{L1}
\difficulte{}

\contenu{
\question{
Soit $(u_{n})_{n\in \Nn}$ la suite d\'{e}finie par $u_{n+1}=u_{n}+u_{n}^{2}. $
L'\'{e}tudier et, en utilisant $v_{n}=\frac{1}{u_{n}},$ en donner un
\'{e}quivalent dans le cas $u_{0}\in ]-1;0]$. Que dire dans le cas $u_{0}\in
]0;\infty [$ ? (On \'{e}tudiera $v_{n}=\frac{\ln (u_{n})}{2^{n}}$.)
}
}
